% SEGMENT OLP-0421-S01
% Part: normal-modal-logic
% Chapter: frame-correspondence
% Section: properties-accessibility

\documentclass[../../../include/open-logic-section]{subfiles}

\begin{document}

\olfileid{nml}{frd}{acc}

\olsection{அணுகுறவுகளின் பண்புகள்}

பல வாய்ப்புநிலை !!{formula}s, அணுகுறவின் எளிய, நன்கு அறிமுகமான பண்புகளுக்குக்
கூட, சிறப்பியல்பாக அமைகின்றன. ஒரு திசையில், ஒரு குறிப்பிட்ட பண்பைக் கொண்ட
எந்த மாதிரியும் அதற்குரிய ஒரு !!a{formula} ஐயும் (அதன் எல்லாப் பதிலீட்டு
நிகழ்வுகளையும்) மெய்யாக்குகிறது என்பது இதன் பொருள். அணுகுறவுகளின் ஐந்து
செவ்வியல் வகைகளையும், அவை மெய்யாக இருப்பதை உறுதிசெய்யும் வாய்பாடுகளையும்
கொண்டு தொடங்குகிறோம்.

\begin{thm}\ollabel{thm:soundschemas}
  $\mModel{M} = \tuple{W, R, V}$ ஒரு மாதிரியாக இருக்கட்டும்.
  \olref{tab:five} இன் இடப்பக்கத்திலுள்ள பண்பு~$R$ க்கு இருந்தால்,
  வலப்பக்கத்திலுள்ள !!a{formula} இன் ஒவ்வொரு நிகழ்வும்~$\mModel{M}$ இல்
  மெய்யாகும்.
\end{thm}

\begin{table}[t]
    \begin{tabular}{| l || l |}
      \hline
      {\emph{$R$ \dots எனில்}} & {\emph{\dots ஆனது~$\mModel{M}$ இல் மெய்யாகும்:}} \\
      \hline \hline
      \emph{தொடர் பண்புள்ளது}: $\forall u \exists v Ruv$ & \hbox to.43\textwidth
           {$\Box p \lif \Diamond p$ \hfill (\Ax{D})} \\
      \hline
      \emph{தற்சுட்டுப் பண்புள்ளது}: $\forall w Rww$
      & \hbox to.43\textwidth
          {$\Box p \lif p$ \hfill (\Ax{T})} \\
      \hline
      \emph{சமச்சீர்ப் பண்புள்ளது}: &  \hbox to .43\textwidth
           {$p \lif \Box\Diamond p$ \hfill (\Ax{B})} \\
           $\forall u\forall v(Ruv \lif Rvu)$ & \\
      \hline
      \emph{கடப்புப் பண்புள்ளது}: & \hbox to.43\textwidth
           {$\Box p \lif \Box \Box p$ \hfill (\Ax{4})} \\
      $\forall u \forall v \forall w ((Ruv \land Rvw) \lif Ruw)$ & \\
      \hline
      \emph{யூக்ளிடியப் பண்புள்ளது}: & \hbox to .43\textwidth
        {$\Diamond p \lif \Box\Diamond p$ \hfill (\Ax{5})} \\
      $\forall w \forall u \forall v ((Rwu \land Rwv) \lif Ruv)$ &\\
      \hline
    \end{tabular}
    \caption{ஐந்து ஒத்திசைவு உண்மைகள்.}
    \ollabel{tab:five}
  \end{table}


\begin{proof}
  \Ax{B} க்கான நேர்வு இதோ: இந்தத் திட்டவடிவம் ஒரு மாதிரியில் மெய்யாகும்
  என்று காட்ட, அதன் எல்லா நிகழ்வுகளும் மாதிரியிலுள்ள எல்லா உலகங்களிலும்
  மெய்யாகும் என்று காட்ட வேண்டும். ஆகவே $!A \lif \Box\Diamond !A$ என்பது
  \Ax{B} இன் தரப்பட்ட ஒரு நிகழ்வாகவும், $w \in W$ ஒரு ஏதேனும் உலகமாகவும்
  இருக்கட்டும். $w$ இல் $\Box \Diamond !A$ மெய்யாகும் என்று காட்டுவதற்காக,
  அதன் முன்னுறுப்பு~$!A$ ஆனது $w$ இல் மெய்யாக உள்ளது என்க. எனவே $w$
  இலிருந்து அணுகத்தக்க எல்லா $w'$ களிலும் $\Diamond !A$ மெய்யாகும் என்று
  காட்ட வேண்டும். இப்போது, $Rww'$ ஆகும் எந்த $w'$ க்கும், சமச்சீர்ப் பண்பு
  என்ற கருதுகோளைப் பயன்படுத்தி $Rw'w$ உம் கிடைக்கிறது (காண்க
  \olref{fig:Bsymm}). $\mSat{M}{!A}[w]$ என்பதால்,
  $\mSat{M}{\Diamond !A}[w']$. $Rww'$ ஆகும் ஏதேனும் உலகமாக $w'$ இருந்ததால்,
  $\mSat{M}{\Box\Diamond!A}[w]$.

  மற்ற நேர்வுகளைப் பயிற்சிகளாக விடுகிறோம்.
\end{proof}

\begin{prob}
  \olref[nml][frd][acc]{thm:soundschemas} இன் நிறுவலை நிறைவுசெய்க.
\end{prob}

\begin{figure}
  \begin{center}
    \begin{tikzpicture}[modal]
      \node[world] (w1) [label={below:$\mSat{{}}{\formula{A}}$\\
          $\mSat{{}}{\Box\Diamond \formula{A}}$}] {$w$};
      \node[world] (w2) [label=below:{$\mSat{{}}{\Diamond \formula{A}}$},
        right=of w1] {$w'$};
      \draw[->, bend left] (w1) to (w2);
      \draw[->, bend left] (w2) to (w1);
    \end{tikzpicture}
  \end{center}
\caption{சமச்சீர்ப் பண்பிலிருந்து பெறும் வாதம்.}
\ollabel{fig:Bsymm}
\end{figure}

\olref{thm:soundschemas} இன் மறுதலை உட்கிடைகள் செல்லாது என்பதைக்
கவனிக்க: ஒரு மாதிரி ஒரு திட்டவடிவத்தை மெய்யாக்கினால், அந்த மாதிரியின்
அணுகுறவிற்கு அதற்குரிய பண்பு இருக்கும் என்பது உண்மையல்ல. \Ax{T} மற்றும்
தற்சுட்டு மாதிரிகளின் நேர்வில், \Ax{T} தானே தவறும் ஒரு மாதிரிக்கான
எடுத்துக்காட்டைக் கொடுப்பது எளிது: $W = \{w\}$ என்றும் $V(p) = \emptyset$
என்றும் கொள்க. அப்போது $R$ தற்சுட்டுப் பண்புடையதல்ல; ஆனால்
$\mSat{M}{\Box p}[w]$ மற்றும் $\mSat/{M}{p}[w]$. எனினும், இங்கு
\Ax{T} இன் ஓர் நிகழ்வு மட்டுமே~$\mModel{M}$ இல் தவறுகிறது; எடுத்துக்காட்டாக,
$\Box \lnot p \lif \lnot p$ போன்ற மற்ற நிகழ்வுகள் மெய்யாகின்றன.
\Ax{T} இன் \emph{ஒவ்வொரு பதிலீட்டு நிகழ்வும்}~$\mModel{M}$ இல் மெய்யாகவும்,
$\mModel{M}$ தற்சுட்டுப் பண்பற்றதாகவும் இருக்கும் எடுத்துக்காட்டுகளைக்
கொடுப்பது கடினம். ஆனால் அத்தகைய மாதிரிகளும் உள்ளன:

\begin{prop}\ollabel{prop:reflexive}
  $\mModel{M} = \tuple{W, R, V}$ என்பது $W = \{u, v\}$ ஆகும் ஒரு மாதிரியாக
  இருக்கட்டும்; இங்கு உலகங்கள் $u$, $v$ ஆகியவை~$R$ ஆல் தொடர்புபடுத்தப்பட்டுள்ளன:
  அதாவது $Ruv$, $Rvu$ இரண்டும் உள்ளன. எல்லா $p$ க்கும் $u \in V(p)
  \Leftrightarrow v \in V(p)$ என்க. அப்போது:
  \begin{enumerate}
  \item எல்லா $!A$ க்கும், $\mSat{M}{!A}[u]$ எனில், எனில் மட்டுமே,
    $\mSat{M}{!A}[v]$ ($!A$ மீதான தொகுத்தறிதலைப் பயன்படுத்துக).
  \item \Ax{T} இன் ஒவ்வொரு நிகழ்வும்~$\mModel{M}$ இல் மெய்யாகும்.
  \end{enumerate}
  $\mModel{M}$ தற்சுட்டுப் பண்புடையதல்ல (உண்மையில் அது
  \emph{எங்கும் தற்சுட்டற்றது}) என்பதால், \Ax{T} இன் நேர்வில்
  \olref{thm:soundschemas} இன் மறுதலை தவறுகிறது (\olref{thm:soundschemas}
  இல் குறிப்பிடப்பட்டுள்ள மற்ற திட்டவடிவங்களில் சிலவற்றுக்கும்---அனைத்திற்கும்
  அல்ல---இதே போன்ற வாதங்களைக் கொடுக்கலாம்).
\end{prop}

\begin{prob}
  \olref[nml][frd][acc]{prop:reflexive} இலுள்ள கூற்றுகளை நிறுவுக.
\end{prob}

ஐந்து செவ்வியல் !!{formula}s ஆன \Ax{D}, \Ax{T}, \Ax{B}, \Ax{4},
\Ax{5} ஆகியவற்றில் நாம் கவனம் செலுத்தப்போகிறோம் என்றாலும், அணுகுறவுகளின்
மேலும் சில பண்புகளை \olref{tab:anotherfive} இல் பதிவு செய்கிறோம்.
ஒவ்வோர் உலகிலிருந்தும் அதிகபட்சம் ஓர் உலகை மட்டுமே அணுக முடிந்தால்,
அணுகுறவு~$R$ பகுதிச் சார்பாக உள்ளது. ஒவ்வோர் உலகிலிருந்தும் சரியாக
ஓர் உலகை அணுக முடிந்தால், அதைச் சார்பாக உள்ளது என்கிறோம். (ஆகவே,
தொடர் பண்பும் பகுதிச் சார்புப் பண்பும் இரண்டையும் கொண்டவையே துல்லியமாகச்
சார்புகளாகும்.) அணுகுறவு ஒரு (பகுதிச்) சார்பைப் போலச் செயல்படுவதால் அவை
``சார்புகள்'' எனப்படுகின்றன. $Ruv$ ஆகும்போதெல்லாம் $u$, $v$ ஆகியவற்றுக்கு
``இடையில்'' ஒரு~$w$ இருந்தால், ஒரு தொடர்பு தளர்வாக அடர்ந்தது.
எனவே தளர்வாக அடர்ந்த தொடர்புகள் ஒரு பொருளில் கடப்புத் தொடர்புகளுக்கு
எதிரானவை: கடப்புத் தொடர்பில்,~$w$ வழியான மாற்றுப்பாதையில் $u$ இலிருந்து
$v$ ஐ அடைய முடியும்போதெல்லாம், $u$ இலிருந்து $v$ ஐ நேரடியாக அடையலாம்;
தளர்வாக அடர்ந்த தொடர்பில், $u$ இலிருந்து $v$ ஐ நேரடியாக அடைய
முடியும்போதெல்லாம், ஏதோ~$w$ வழியான மாற்றுப்பாதையிலும் அதை அடையலாம்.
ஏதோ ஓர் உலகம்~$w$ இலிருந்து $u$, $v$ ஆகிய உலகங்களை அடைய
முடியும்போதெல்லாம், $u$, $v$ இரண்டிலிருந்தும் ஒரு பொது உலகம்~$t$ ஐ அடைய
முடிந்தால், தொடர்பு தளர்வாகத் திசைப்படுத்தப்பட்டது; இது சில நேரங்களில்
``வைரப் பண்பு'' அல்லது ``சங்கமம்'' எனப்படுகிறது.

\begin{table}[t]
    \begin{tabular}{| p{.50\textwidth} || p{.43\textwidth} |}
      \hline
      {\emph{$R$ \dots எனில்}} & {\emph{\dots ஆனது~$\mModel{M}$ இல் மெய்யாகும்:}} \\
      \hline \hline
      \emph{பகுதிச் சார்பாக உள்ளது}: &
      \multirow{2}{*}{$\Diamond p \lif \Box p$} \\
      $ \forall w \forall u \forall v ((Rwu \land Rwv) \lif u =v )$ & \\
      \hline
      \multirow{2}{*}{\emph{சார்பாக உள்ளது}: $\forall w \exists v \forall u(Rwu \liff \eq[u][v]) $}
      & \multirow{2}{*}{$\Diamond p \liff \Box p$} \\
      & \\
      \hline
      \emph{தளர்வாக அடர்ந்தது}: &  \multirow{2}{*}{$\Box \Box p \lif
        \Box p$} \\
      $\forall u\forall v(Ruv \lif \exists w(Ruw \land Rwv))$ & \\
      \hline
      \emph{தளர்வாக இணைந்தது}: &
      \multirow{3}{*}{\hbox to.43\textwidth{$
        \begin{array}{@{}l@{}}
          \Box ((p \land \Box p) \lif q) \lor {} \\
          \qquad \Box ((q \land \Box q) \lif p)
        \end{array}$ \hfill (\Ax{L})}
      } \\
      $\forall w \forall u \forall v ((Rwu \land Rwv) \lif {}$ &\\
      \qquad $(Ruv \lor u=v \lor Rvu))$ & \\
      \hline
      \emph{தளர்வாகத் திசைப்படுத்தப்பட்டது}: & \multirow{3}{*}{\hbox to .43\textwidth
        {$\Diamond\Box p \lif \Box\Diamond p$\hfill (\Ax{G})}} \\
      $\forall w \forall u \forall v ((Rwu \land Rwv) \lif {}$ & \\
      \qquad $\exists t (Rut \land Rvt))$ & \\
      \hline
    \end{tabular}
    \caption{மேலும் ஐந்து ஒத்திசைவு உண்மைகள்.}
    \ollabel{tab:anotherfive}
  \end{table}

\begin{prob}
  $\mModel{M} = \tuple{W, R, V}$ ஒரு மாதிரியாக இருக்கட்டும். $R$ ஆனது
  \olref[nml][frd][acc]{tab:anotherfive} இன் இடப்பக்கப் பண்புகளை
  நிறைவுசெய்தால், அதற்குரிய வலப்பக்க வாய்பாட்டின் ஒவ்வொரு நிகழ்வும்
  $\mModel{M}$ இல் மெய்யாகும் என்று காட்டுக.
\end{prob}



\end{document}
